\documentclass{article}

\usepackage{parskip}
\usepackage{hyperref}
\usepackage{array}
\usepackage{enumitem}
\usepackage{longtable}
\usepackage{listings}
\usepackage{xcolor}

\hypersetup{
  pdftitle={Java Code Conventions},
  pdfauthor={Larry Kramer},
  pdfsubject={Java Code Conventions},
  pdfkeywords={Java, code conventions, coding conventions, coding standards, style guide, best practices},
  hidelinks,
  linktoc=all
}
\lstset{
  aboveskip=0.6\baselineskip,
  belowskip=0.6\baselineskip,
  basicstyle=\ttfamily\footnotesize,      % shrink just enough to fit
  tabsize=4,
  showstringspaces=false,                 % hide whitespace markers
  xleftmargin=0pt,                        % ensure the same left edge on every page
  xrightmargin=0pt,
  breaklines=false,                       % NO line breaks
  breakatwhitespace=false,
  framexleftmargin=0pt,
  columns=fullflexible,                   % allow tighter glyph boxes in \ttfamily
  keepspaces=true                         % preserve indentation alignment
}

\newlist{notes}{enumerate}{1}
\setlist[notes]{label=\textbf{Note:}, leftmargin=*}
\newenvironment{note}[1]
  {\begin{notes}\item #1}{\end{notes}}

\lstnewenvironment{snippet}[1][]
  {\noindent\minipage{\linewidth}}
  {\endminipage}

\begin{document}

\begin{titlepage}
  \centering
  \vspace*{0.33\textheight}
  {\Huge\bfseries Java\par}
  \vspace{0.5em}
  {\Huge\bfseries Code Conventions\par}
  \vspace{1.5em}
  {\large\bfseries \today\par}
  \vfill
\end{titlepage}

\tableofcontents
\clearpage

\section{Introduction}

This document outlines the standard Java coding conventions for all software development projects.
The primary goal of these standards is to improve code readability, consistency, and long-term
maintainability across the entire codebase.

These conventions should be treated as guidelines, not as immutable rules. The ultimate goal is
always clarity and maintainability. Developers are expected to apply these conventions judiciously.
In specific contexts where adhering to a guideline would degrade code quality, developers should
prioritize the clearer solution and be prepared to justify their deviations during code review.

Adhering to these conventions ensures that all developers can easily understand, review, and
contribute to the project.

\subsection{Acknowledgments}

These conventions are largely based on the
\href{https://www.oracle.com/docs/tech/java/codeconventions.pdf}{Sun/Oracle
\emph{Java Code Conventions} (1997 Edition)}, with specific additions and modernizations tailored to
the demands of contemporary software engineering.

\section{File Names}

This section lists commonly used file suffixes and names.

\subsection{File Suffixes}

Frequently used file suffixes include:

\begin{tabular}{ll}
\hline
File Type & Suffix \\
\hline
Java source & \texttt{.java} \\
Java class file & \texttt{.class} \\
XML document & \texttt{.xml} \\
Markdown document & \texttt{.md} \\
\hline
\end{tabular}

\subsection{Common File Names}

Frequently used file names include:

\begin{tabular}{p{0.2\linewidth} >{\raggedright\arraybackslash}p{0.8\linewidth}}
\hline
File Name & Use \\
\hline
\texttt{pom.xml} & The preferred name for the build file. \newline We use Apache Maven to build our
  software. \\
\texttt{README.md} & The preferred name for the file that summarizes the contents of a particular
  directory. \\
\hline
\end{tabular}

\section{File Organization}

A file consists of several sections, separated by blank lines. Optionally, a comment can be used to
identify each section.

Java source files longer than 2,000 lines are cumbersome and should be avoided.

For an example of a properly formatted Java program, see
\hyperref[codeexample]{"\nameref*{codeexample}" on page~\pageref*{codeexample}}.

\subsection{Java Source Files}

Each Java source file contains a single public class or interface. When non-public classes and
interfaces are associated with a public class, they can be placed in the same source file. The
public class should be the first class or interface in the file.

Java source files have the following ordering:

\begin{enumerate}
\item Beginning comments (see \hyperref[beginningcomments]{"\nameref*{beginningcomments}" on
  page~\pageref*{beginningcomments}})
\item Package and import statements (see \hyperref[packagesimports]{"\nameref*{packagesimports}" on
  page~\pageref*{packagesimports}}); for example: \\
\begin{snippet}
import java.net.*;
import java.util.*;

import net.larrykramer.somepackage.SomeClass;
\end{snippet}
\item Class and interface declarations (see
  \hyperref[classesinterfaces]{"\nameref*{classesinterfaces}" on page~\pageref*{classesinterfaces}})
\end{enumerate}

\subsubsection{Beginning Comments}
\label{beginningcomments}

All source files should begin with a C-style comment that lists a copyright notice. For example:

\begin{snippet}
/*
 * Copyright notice
 */
\end{snippet}

\subsubsection{Package and Import Statements}
\label{packagesimports}

The first non-comment line of most Java source files is a package statement. After the package
statement, \texttt{import} statements can follow.

Imports must be ordered in the following groups, with a blank line separating each group. Within
each group, imports must be sorted alphabetically (ASCII order).

\begin{enumerate}
\item All module imports
\item \texttt{java.*} and \texttt{javax.*} libraries
\item All other non-static imports (third-party, project-specific)
\item All static imports
\end{enumerate}

Use a wildcard import (e.g., \texttt{import java.util.*;}) only when \emph{four or more} distinct
types are imported from the exact same package. The "four or more" rule also applies to the static
members of a single, fully-qualified class. This rule applies to both standard (non-static) and
static imports, but the groups are counted independently.

\begin{snippet}
// Group 1: Module import declarations
import module java.logging;

// Group 2: Standard Java libraries
import java.io.File;
import java.util.List;

// Group 3: Other non-static imports
import org.junit.Test;
import org.mockito.*;
import org.mockito.junit.MockitoJUnitRunner;

// Group 4: All static imports
import static java.lang.Math.max;
import static org.junit.Assert.*;
import static org.mockito.Mockito.*;
\end{snippet}

\subsubsection{Class and Interface Declarations}
\label{classesinterfaces}

The following table describes the parts of a class or interface declaration, in the order that they
should appear. See \hyperref[codeexample]{"\nameref*{codeexample}" on page~\pageref*{codeexample}}
for an example that includes comments.

\begingroup
\renewcommand{\arraystretch}{1.5}
\begin{longtable}{
  p{0.01\linewidth}|
  >{\raggedright\arraybackslash}p{0.37\linewidth}
  >{\raggedright\arraybackslash}p{0.52\linewidth}}
\hline
\# & Part of Class/Interface & Declaration Notes \\
\hline
\endhead
\hline
\endlastfoot
{1} & Class/interface documentation comment (\texttt{/**...*/}) & See
  \hyperref[doccomments]{"\nameref*{doccomments}" on page~\pageref*{doccomments}}
  for information on what this comment should contain. \\
2 & \texttt{class}, \texttt{interface}, \texttt{enum}, or \texttt{record} statement & \\
3 & Class/interface implementation comment (\texttt{/*...*/}), if necessary & This comment should
  contain any class-wide or interface-wide information that was not appropriate for the
  class/interface documentation comment. \\
4 & Class (\texttt{static}) variables & First the \texttt{public} class variables, then
  \texttt{protected}, and then \texttt{private}. \\
5 & Instance variables & First \texttt{public}, then \texttt{protected}, and then
  \texttt{private}. \\
6 & Constructors & \\
7 & Methods & These methods should be grouped by functionality rather than by scope or
  accessibility. For example, a private class method can be placed between two public instance
  methods. The goal is to make reading and understanding the code easier. \\
\end{longtable}
\endgroup

\section{Indentation}

Four spaces should be used as the unit of indentation. Spaces are the only characters permitted
for indentation; the use of tab characters is forbidden.

\begin{note}
Some tools may require tab characters for alignment within string literals or comments. This is a
rare exception and does not apply to code indentation. If tabs are used in these specific cases,
they must be configured to be equivalent to 8 spaces.
\end{note}

\subsection{Line Length}

Avoid lines longer than 100 characters, since they are not handled well by many terminals and
tools.

\begin{note}
Examples for use in documentation should have a shorter line length---generally no more than 90
characters.
\end{note}

\subsection{Wrapping Lines}

When an expression will not fit on a single line, break it according to these general principles:

\begin{itemize}
\item Break after a comma.
\item Break before an operator.
\item Prefer higher-level breaks to lower-level breaks.
\item Align the new line with the beginning of the expression at the same level on the previous
  line.
\item If the above rules lead to confusing code or to code that is squished up against the right
  margin, just indent 8 spaces instead.
\end{itemize}

Here are some examples of breaking method calls:

\begin{snippet}
function(longExpression1, longExpression2, longExpression3,
        longExpression4, longExpression5);

aVariable = function1(longExpression1,
        function2(longExpression2,
                longExpression3));
\end{snippet}


The following are two examples of breaking an arithmetic expression. The first is preferred, since
the break occurs outside the parenthesized expression, which is at a higher level.

\begin{snippet}
longName1 = longName2 * (longName3 + longName4 - longName5)
            + 4 * longName6; // PREFERRED

longName1 = longName2 * (longName3 + longName4
                         - longName5) + 4 * longName6; // AVOID!
\end{snippet}

The following is an example of indenting method declarations.

\begin{snippet}
// CONVENTIONAL INDENTATION
someMethod(int anArg, Object anotherArg, String yetAnotherArg,
        Object andStillAnother) {
 ...
}
\end{snippet}

Line wrapping for if statements should generally use the 8-space rule, since conventional (4-space)
indentation can obscure the body. For example:

\begin{snippet}
// DO NOT USE THIS INDENTATION
if ((condition1 && condition2)
    || (condition3 && condition4)
    || !(condition5 && condition6)) { // BAD WRAPS
    doSomethingAboutIt(); // MAKES THIS LINE EASY TO MISS
}

// USE THIS INDENTATION INSTEAD
if ((condition1 && condition2)
        || (condition3 && condition4)
        || !(condition5 && condition6)) {
    doSomethingAboutIt();
}

// OR USE THIS
if ((condition1 && condition2) || (condition3 && condition4)
        || !(condition5 && condition6)) {
    doSomethingAboutIt();
}
\end{snippet}

Here are two acceptable ways to format ternary expressions:

\begin{snippet}
alpha = (aLongBooleanExpression) ? beta : gamma;

alpha = (aLongBooleanExpression)
        ? beta
        : gamma;
\end{snippet}

\section{Comments}

Java programs can have two kinds of comments: implementation comments and documentation comments.
Implementation comments use the same syntax as C++, delimited by \texttt{/*...*/} and
\texttt{//}. Documentation comments (known as "doc comments") are specific to Java and are delimited
by \texttt{/**...*/}. Doc comments can be extracted to HTML files using the Javadoc tool.

Implementation comments are meant for commenting out code or for comments about the particular
implementation. Doc comments are meant to describe the specification of the code, from an
implementation-free perspective, to be read by developers who might not necessarily have the source
code at hand.

Comments should be used to give an overview of the code and provide additional information that is
not readily available in the code itself. Comments should contain only information that is relevant
to reading and understanding the program. For example, information about how the corresponding
package is built or in what directory it resides should not be included as a comment.

Discussion of non-trivial or non-obvious design decisions is appropriate, but avoid duplicating
information that is present in (and clear from) the code. It is too easy for redundant comments to
become outdated. In general, avoid any comments that are likely to get out of date as the code
evolves.

\begin{note}
The frequency of comments sometimes reflects poor quality code. When you feel compelled to add a
comment, consider rewriting the code to make it clearer.
\end{note}

Comments should not be enclosed in large boxes drawn with asterisks or other characters. Comments
should never include special characters such as form-feed and backspace.

\subsection{Implementation Comment Formats}

Programs can have four styles of implementation comments: block, single-line, trailing, and
end-of-line.

\subsubsection{Block Comments}
\label{blockcomments}

Block comments are used to provide descriptions of files, methods, data structures, and algorithms.
Block comments should be used at the beginning of each file and before each method. They can also be
used in other places, such as within methods. Block comments inside a function or method should be
indented to the same level as the code they describe.

Block comments should have an asterisk "\texttt{*}" at the beginning of each line except the first.

\begin{snippet}
/*
 * Here is a block comment.
 */
\end{snippet}

\subsubsection{Single-Line Comments}
\label{singlelinecomments}

Short comments can appear on a single line indented to the level of the code that follows. If a
comment cannot be written in a single line, it should follow the block comment format (see
\hyperref[blockcomments]{section~\ref*{blockcomments} on page~\pageref*{blockcomments}}). Here is
an example of a single-line comment in Java code (also see
\hyperref[doccomments]{"\nameref*{doccomments}" on page~\pageref*{doccomments}}):

\begin{snippet}
if (condition) {
    /* Handle the condition. */
    ...
}
\end{snippet}

\subsubsection{Trailing Comments}

Very short comments can appear on the same line as the code they describe, but should be shifted far
enough to separate them from the statement. If more than one short comment appears in a chunk of
code, they should all be indented to the same column. Avoid the assembly-language style of placing a
trailing comment on every line of executable code. Here is an example of a trailing comment in Java
code (also see \hyperref[doccomments]{"\nameref*{doccomments}" on page~\pageref*{doccomments}}):

\begin{snippet}
if (a == 2) {
    return true; /* special case */
} else {
    return isPrime(a); /* works only for odd a */
}
\end{snippet}

\subsubsection{End-Of-Line Comments}

The \texttt{//} comment delimiter begins a comment that continues to the newline. It can comment out
a complete line or only a partial line. It should not be used on consecutive multiple lines for text
comments; however, it can be used in consecutive multiple lines for commenting out sections of
code. Examples of all three styles follow:

\begin{snippet}
if (foo > 1) {
    // Do a double-flip.
    ...
} else {
    return false; // Explain why here.
}
//if (bar > 1) {
//
//    // Do a triple-flip.
//    ...
//} else {
//    return false;
//}
\end{snippet}

\subsection{Documentation Comments}
\label{doccomments}

\begin{note}
See \hyperref[codeexample]{"\nameref*{codeexample}" on page~\pageref*{codeexample}} for examples of
the comment formats described here.
\end{note}

For further details, see "How to Write Doc Comments for Javadoc", which includes information on the
doc comment tags (\texttt{@return}, \texttt{@param}, \texttt{@see}):

\begin{quote}
\mbox{\small\url{https://www.oracle.com/technical-resources/articles/java/javadoc-tool.html}\par}
\end{quote}

Doc comments describe Java classes, interfaces, constructors, methods, and fields. Each doc comment
is set inside the comment delimiters \texttt{/**...*/}, with one comment per API. This comment
should appear just before the declaration:

\begin{snippet}
/**
 * The Example class provides ...
 */
class Example { ...
\end{snippet}

Note that the declaration lines for classes and interfaces are not indented, while their members
are. The first line of a doc comment (\texttt{/**}) for classes and interfaces is not indented;
subsequent doc comment lines each have 1 space of indentation (to vertically align the
asterisks). Members, including constructors, have 4 spaces for the first doc comment line and 5
spaces thereafter.

If you need to give information about a class, interface, variable, or method that is not
appropriate for documentation, use an implementation block comment (see
\hyperref[blockcomments]{section~\ref*{blockcomments} on page~\pageref*{blockcomments}}) or
single-line (see \hyperref[singlelinecomments]{section~\ref*{singlelinecomments} on
page~\pageref*{singlelinecomments}} comment immediately \emph{after} the declaration. For example,
details about the implementation of a class belong in an implementation block comment following the
class statement, not in the class doc comment.

Doc comments should not be positioned inside a method or constructor definition block, because Java
associates documentation comments with the first declaration \emph{after} the comment.

\section{Declarations}

\subsection{Number Per Line}

One declaration per line is recommended since it encourages commenting. In other words,

\begin{snippet}
int level; // indentation level
int size;  // size of table
\end{snippet}

is preferred over

\begin{snippet}
int level, size;
\end{snippet}

Do not put different types on the same line. Example:

\begin{snippet}
int foo, fooarray[]; // WRONG!
\end{snippet}

\begin{note}
The examples above use one space between the type and the identifier. As an acceptable alternative,
you can use additional spaces to align the variable names and comments.

\begin{snippet}
int    level;        // indentation level
int    size;         // size of table
Object currentEntry; // currently selected table entry
\end{snippet}
\end{note}

\subsection{Placement}

Place declarations as close to their first use as possible. Declarations at the beginning of a
block (a block is any code surrounded by curly braces "\texttt{\{}" and "\texttt{\}}") are also
acceptable.

\begin{snippet}
void myMethod() {
    int int1;            // beginning of method block
    if (condition) {
        ...
        int int2 = ...;  // declared on first use
        ...
    }
}
\end{snippet}

The one exception to the rule is indices of a \texttt{for} loop, which can be declared within the
\texttt{for} statement itself.

\begin{snippet}
for (int i = 0; i < maxLoops; i++) { ...
\end{snippet}

Avoid local declarations that hide declarations at higher levels. For example, do not declare the
same variable name in an inner block:

\begin{snippet}
int count;
...
func() {
    if (condition) {
        int count = ...; // AVOID!
        ...
    }
    ...
}
\end{snippet}

\subsection{Initialization}

Try to initialize local variables where they are declared. The only reason not to initialize a
variable where it is declared is if the initial value depends on some computation occurring first.

\subsection{Local-Variable Type Inference (\texttt{var})}

Avoid using local-variable type inference (\texttt{var}) except in specific situations where the
initializer makes the type unambiguous and the explicit type would be excessively verbose. The
explicit type is preferred in most cases because \texttt{var} can obscure the code's clarity.

Do not use \texttt{var} where the type is not obvious from the right-hand side of the expression.

\begin{snippet}
var activeConnections = new ArrayList<>();          // VERY WRONG! INFERRED AS ArrayList<Object>.

var activeConnections = new ArrayList<String>();    // AVOID!

List<String> activeConnections = new ArrayList<>(); // OK
\end{snippet}

Using \texttt{var} is appropriate when an explicit type declaration is excessively long.

\begin{snippet}
// AVOID! REDUNDANT, VERBOSE, AND REPEATS RHS TYPE INFORMATION!
Set<Map.Entry<String, CopyOnWriteArrayList<Consumer<Object>>>> entries = map.entrySet();

var entries = map.entrySet(); // OK
\end{snippet}

It is also acceptable to use \texttt{var} in enhanced \texttt{for} loops and
\texttt{try-with-resources} statements.

\begin{snippet}
for (var user : users) {
    ...
}

try (var reader = new BufferedReader(new FileReader("file.txt"))) {
    ...
}
\end{snippet}

\subsection{Class and Interface Declarations}

When coding Java classes and interfaces, the following formatting rules should be followed:

\begin{itemize}
\item No space between a method name and the parenthesis "\texttt{(}" starting its parameter list.
\item Open brace "\texttt{\{}" appears at the end of the same line as the declaration statement.
\item Closing brace "\texttt{\}}" starts a line by itself, indented to match its corresponding
  opening statement. \\
\begin{snippet}
class Sample extends Object {
    int ivar1;
    int ivar2;

    Sample(int i, int j) {
        ivar1 = i;
        ivar2 = j;
    }

    void emptyMethod() {
    }

    ...
}
\end{snippet}
\item Methods are separated by a blank line.
\end{itemize}

\section{Statements}

\subsection{Simple Statements}

Each line should contain at most one statement. Example:

\begin{snippet}
argv++; argc--; // AVOID!
\end{snippet}

\subsection{Compound Statements}
\label{compoundstatements}

Compound statements are statements that contain lists of statements enclosed in braces
"\texttt{\{ statements \}}". See the following sections for examples.

\begin{itemize}
\item The enclosed statements should be indented one more level than the compound statement.
\item The opening brace should be at the end of the line that begins the compound statement; the
  closing brace should begin a line and be indented to the beginning of the compound statement.
\item Braces are used around all statements, even singletons, when they are part of a control
  structure, such as an \texttt{if-else} or \texttt{for} statement. This makes it easier to add
  statements without accidentally introducing bugs due to forgetting to add braces.
\end{itemize}

\subsection{\texttt{return} Statements}

A \texttt{return} statement with a value should not use parentheses unless they make the return
value more obvious in some way. Example:

\begin{snippet}
return;

return myDisk.size();

return (Disk) (isHDD ? new HardDiskDrive() : new FloppyDrive());
\end{snippet}

\subsection{\texttt{if}, \texttt{if-else}, \texttt{if-else-if-else} Statements}

The \texttt{if-else} class of statements should have the following form.

\begin{snippet}
if (condition) {
    statements;
}

if (condition) {
    statements;
} else {
    statements;
}

if (condition) {
    statements;
} else if (condition) {
    statements;
} else if (condition) {
    statements;
} else {
    statements;
}
\end{snippet}

\begin{note}
\texttt{if} statements always use braces "\texttt{\{\}}". Avoid the following error-prone form:

\begin{snippet}
if (condition)  // AVOID! THIS OMITS THE BRACES {}!
    statement;
\end{snippet}
\end{note}

\subsection{\texttt{for} Statements}

A \texttt{for} statement should have the following form:

\begin{snippet}
for (initialization; condition; update) {
    statements;
}
\end{snippet}

An empty \texttt{for} statement (one in which all the work is done in the initialization, condition,
and update clauses) can have the following form:

\begin{snippet}
for (initialization; condition; update);
\end{snippet}

When using comma-separated initializers or updaters in a \texttt{for} statement, avoid the
complexity of using more than three variables. If needed, use separate statements before the
\texttt{for} loop (for the initialization clause) or at the end of the loop (for the update clause).

\subsection{\texttt{while} Statements}

A \texttt{while} statement should have the following form:

\begin{snippet}
while (condition) {
    statements;
}
\end{snippet}

An empty \texttt{while} statement can have the following form:

\begin{snippet}
while (condition);
\end{snippet}

\subsection{\texttt{do-while} Statements}

A \texttt{do-while} statement should have the following form:

\begin{snippet}
do {
    statements;
} while (condition);
\end{snippet}

\subsection{\texttt{switch} Statements}

A \texttt{switch} statement should have the following \emph{preferred modern form}, with the arrow
syntax (\texttt{->}), to avoid fall-through logic and improve readability.

\begin{snippet}
switch (condition) {
    case ABC -> statement;
    case DEF -> statement;
    case XYZ -> statement;
    default -> statement;
}
\end{snippet}

A \texttt{switch} statement can have the following traditional form:

\begin{snippet}
switch (condition) {
    case ABC:
        statements;
        /* falls through */
    case DEF:
        statements;
        break;
    case XYZ:
        statements;
        break;
    default:
        statements;
        break;
}
\end{snippet}

Every time a case falls through (i.e., does not include a \texttt{break} statement), add a comment
where the \texttt{break} statement would normally appear. This is shown in the preceding code
example with the \texttt{/* falls through */} comment.

Every traditional \texttt{switch} statement should include a default case. The \texttt{break} in the
default case is redundant, but it prevents a fall-through error if another case is later added.

\subsection{\texttt{try-catch}, \texttt{try-with-resources} Statements}

A \texttt{try-catch} statement should have the following format:

\begin{snippet}
try {
    statements;
} catch (ExceptionClass e) {
    statements;
} finally {
    statements;
}
\end{snippet}

A \texttt{try-with-resources} statement should have the following format:

\begin{snippet}
try (resource-specification) {
    statements;
} catch (ExceptionClass e) {
    statements;
} finally {
    statements;
}
\end{snippet}

A \texttt{try-catch} statement in which the exception object is not used should use an unnamed
variable (\texttt{\_}) to make the intent clear.

\begin{snippet}
try {
    statements;
} catch (ExceptionClass _) {
    // The exception is intentionally ignored.
}
\end{snippet}

\section{White Space}

\subsection{Blank Lines}

Blank lines improve readability by setting off sections of code that are logically related.

Two blank lines should always be used in the following circumstances:

\begin{itemize}
\item Between top-level type declarations (e.g., classes, interfaces, enums) within the same source
  file.
\end{itemize}

One blank line should always be used in the following circumstances:

\begin{itemize}
\item Between sections of a source file.
\item Between methods.
\item Between the local variables in a method and its first statement.
\item Before a block (see \hyperref[blockcomments]{section~\ref*{blockcomments} on
  page~\pageref*{blockcomments}}) or single-line (see
  \hyperref[singlelinecomments]{section~\ref*{singlelinecomments} on
  page~\pageref*{singlelinecomments}}) comment except before the first statement of a compound
  statement (see \hyperref[compoundstatements]{section~\ref*{compoundstatements} on
  page~\pageref*{compoundstatements}}).
\item Between logical sections inside a method to improve readability.
\end{itemize}

\subsection{Blank Spaces}

Blank spaces should be used in the following circumstances:

\begin{itemize}
\item A keyword followed by a parenthesis should be separated by a space. Example: \\
\begin{snippet}
while (true) {
    ...
}
\end{snippet}

  Note that a blank space should not be used between a method name and its opening parenthesis.
  This helps to distinguish keywords from method calls.
%
\item A blank space should appear after commas in argument lists.
%
\item All binary operators except for "\texttt{.}" and "\texttt{::}" should be separated from their
  operands by spaces. Blank spaces should never separate unary operators such as unary minus,
  increment ("\texttt{++}"), and decrement ("\texttt{--}") from their operands.
  Example: \\
\begin{snippet}
a += c + d;
a = (a + b) / (c * d);

while ((d = s++) != 0) {
    n++;
}
System.out.print("size is " + foo + "\n");
\end{snippet}
%
\item The arrow operator (\texttt{->}) in lambda and switch expressions should be surrounded by
  blank spaces. Example: \\
\begin{snippet}
  s -> s.length()

  case "FOO" -> statement;
\end{snippet}
%
\item The expressions in a \texttt{for} statement should be separated by blank spaces. Example: \\
\begin{snippet}
for (expr1; expr2; expr3)
\end{snippet}
%
\item Casts should be followed by a blank space. Example: \\
\begin{snippet}
myMethod((byte) aNum, (Object) x);
myFunc((int) (cp + 5), ((int) (i + 3)) + 1);
\end{snippet}
\end{itemize}

\section{Naming Conventions}

Naming conventions make programs more understandable by making them easier to read.

They can also convey information about the function of the identifier---for example, whether it is a
constant, package, or class---which can be helpful in understanding the code.

The conventions given in this section are high-level.

\begingroup
\renewcommand{\arraystretch}{1.5}
\begin{longtable}{
  >{\raggedright\arraybackslash}p{0.125\linewidth}|
  >{\raggedright\arraybackslash}p{0.4\linewidth}
  >{\raggedright\arraybackslash}p{0.375\linewidth}}
\hline
Identifier Type & Rules for Naming & Example \\
\hline
\endhead
\hline
\endlastfoot
%
Classes & Class names should be nouns, in mixed case with the first letter of each internal word
  capitalized. Try to keep your class names simple and descriptive. Use whole words—avoid acronyms
  and abbreviations (unless the abbreviation is much more widely used than the long form, such as
  URL or HTML). &
  \lstinline`class NetworkManager;` \newline
  \lstinline`class Session;` \newline
  \lstinline`class EventQueue;` \\
%
Interfaces & Interface names should be capitalized like class names. &
  \lstinline`interface Listener` \\
%
Records & Record names should be capitalized like class names. &
   \lstinline`record ConnectionDetail` \\
%
Enums & Enum names should be capitalized like class names. &
   \lstinline`enum NetworkStatus` \\
%
Methods & Methods should be verbs, in mixed case with the first letter lowercase, with the first
  letter of each internal word capitalized. &
  \lstinline`openConnection(url);` \newline
  \lstinline`getActiveConnectionCount();` \newline
  \lstinline`logStatus(status);` \\
%
Variables & Instance, class, and local variables are written in mixed case with a lowercase first
  letter. Internal words start with capital letters.
  \par\vspace{0.25\baselineskip plus 2pt}
  Variable names should be short yet meaningful. The choice of a variable name should be
  mnemonic---that is, designed to indicate to the casual observer the intent of its use.
  One-character variable names should be avoided except for temporary "throwaway" variables. Common
  names for temporary variables are \texttt{i}, \texttt{j}, \texttt{k}, \texttt{m}, and \texttt{n}
  for integers; \texttt{c}, \texttt{d}, and \texttt{e} for characters. When a variable is
  syntactically required but intentionally unused, an unnamed variable (\texttt{\_}) should be
  used instead of a name like \texttt{ignored}. &
  \lstinline`int maxConnections;` \newline
  \lstinline`char c;` \newline
  \lstinline`int i;` \\
%
Constants & The names of variables declared as class constants should be all uppercase with words
  separated by underscores ("\_"). &
  \lstinline`static final int MAX_CONNECTIONS = 10;` \newline
  \lstinline`static final int TIMEOUT_FACTOR = 2;` \newline
  \lstinline`static final int MAX_LENGTH = 1024;` \\
\end{longtable}
\endgroup

\section{Programming Practices}

\subsection{Providing Access to Instance and Class Variables}

Do not make any instance or class variable public without good reason. Often, instance variables
do not need to be explicitly set or accessed directly—often that happens as a side effect of method
calls.

One example of appropriate public instance variables is the case where the class is essentially a
mutable data structure, with no behavior. In other words, if you would have used a C-style
\texttt{struct} instead of a class (if Java supported structs), then it is appropriate to make the
class's instance variables public.

For the immutable equivalent, where a class acts as a simple data structure with no behavior, modern
Java provides a more suitable construct. In these cases, you should prefer using a record type,
which is specifically designed to transparently and safely hold immutable data.

\subsection{Referring to Class Variables and Methods}

Avoid using an object to access a class (\texttt{static}) variable or method. Use a class name
instead. For example:

\begin{snippet}
classMethod();           // OK
AClass.classMethod();    // OK
anObject.classMethod();  // AVOID!
\end{snippet}

\subsection{Constants}

Numerical constants (literals) should not be coded directly, except for \texttt{-1}, \texttt{0},
and \texttt{1}, which can appear in a \texttt{for} loop as counter values.

\subsection{Variable Assignments}

Avoid assigning several variables to the same value in a single statement. It is hard to read.
Example:

\begin{snippet}
fooBar.fChar = barFoo.lchar = 'c'; // AVOID!
\end{snippet}

Do not use the assignment operator in a context where it can be easily confused with the equality
operator. Example:

\begin{snippet}
if (aBoolean = isPrime(b)) {  // AVOID! ASSIGNS INSTEAD OF COMPARES!
    ...
}
\end{snippet}

should be written as

\begin{snippet}
if (aBoolean == isPrime(b)) {
    ...
}
\end{snippet}

Do not use embedded assignments in an attempt to improve run-time performance. This is the
job of the compiler, and besides, it rarely actually helps. Example:

\begin{snippet}
d = (a = b + c) + r;  // AVOID!
\end{snippet}

should be written as

\begin{snippet}
a = b + c;
d = a + r;
\end{snippet}

\subsection{Miscellaneous Practices}

\subsubsection{Parentheses}

Avoid using redundant parentheses when operator precedence is clear and the intended order of
operations is unambiguous. While parentheses are essential for controlling expression evaluation,
their overuse can clutter code and impair readability without adding clarity.

Only add parentheses when they are required to enforce a specific order of operations that overrides
the language's default precedence, or in complex expressions where the precedence could be ambiguous
to other programmers. For standard logical and arithmetic expressions, you should rely on the
well-established rules of precedence.

\begin{snippet}
if (a == b && c == d)     // RIGHT

System.out.println("Total: "
        + ((a > 0 || b > 0) ? (a + b) + " dollars" : "N/A"));  // RIGHT

if ((a == b) && (c == d)) // AVOID!
\end{snippet}

\subsubsection{Returning Values}

Try to make the structure of your program match the intent. Example:

\begin{snippet}
if (booleanExpression) {
    return true;
} else {
    return false;
}
\end{snippet}

should instead be written as

\begin{snippet}
return booleanExpression;
\end{snippet}

Similarly,

\begin{snippet}
if (condition) {
    return x;
}
return y;
\end{snippet}

should be written as

\begin{snippet}
return condition ? x : y;
\end{snippet}

\subsubsection{Expressions before ? in the Conditional Operator}

If an expression containing a binary operator appears before the "\texttt{?}" in the ternary
\texttt{?:} operator, it should be parenthesized. Example:

\begin{snippet}
(x >= 0) ? x : -x
\end{snippet}

\subsubsection{Special Comments}

Use \texttt{XXX} in a comment to flag something that is bogus but works. Use \texttt{FIXME} to flag
something that is bogus and broken. Use \texttt{TODO} to flag a planned enhancement or a piece of
missing but non-critical functionality.

\section{Modern Language Features}

This section covers conventions related to features introduced in modern versions of the Java
language.

\subsection{Lambda Expressions}

Lambda expressions provide a concise syntax for representing anonymous functions and are a key
feature of modern Java. Their use is recommended for reducing boilerplate code and improving the
clarity of functional interfaces.

\subsubsection{Single-Expression Lambdas}

A lambda expression whose body consists of a single expression should not be enclosed in braces
\texttt{\{\}}. The \texttt{return} keyword is implicit and must not be specified.

\begin{snippet}
list.stream().map(s -> { return s.toUpperCase(); }); // AVOID!

list.stream().map(s -> s.toUpperCase());             // OK
\end{snippet}

\subsubsection{Lambda Parameters}

The types of lambda parameters should be omitted whenever the compiler can infer them. Avoid
parentheses around a single, inferred-type parameter.

\begin{snippet}
// AVOID
Comparator<String> c = (String s1, String s2) -> s1.compareTo(s2);
Function<String, String> f = (s) -> s.toUpperCase();

// RIGHT
Comparator<String> c = (s1, s2) -> s1.compareTo(s2);
Function<String, String> f = s -> s.toUpperCase();
\end{snippet}

\subsubsection{Complex Lambdas}

While lambdas offer conciseness, their primary goal should be to enhance, not hinder, code clarity.
There is no strict rule for the maximum length of a lambda; many real-world scenarios, particularly
in asynchronous programming with \texttt{CompletableFuture} or in complex data transformations,
involve multi-line lambdas that are perfectly acceptable.

The decision to refactor a lambda into a separate private method should be guided by the following
principles:

\begin{itemize}
\item A lambda should be extracted if its body is so long or complex that it obscures the primary
  purpose of the enclosing method (e.g., the flow of a stream pipeline).
\item A lambda should be extracted if the same block of logic is needed in more than one place, to
  adhere to the DRY (Don't Repeat Yourself) principle.
\item A lambda should be extracted if its internal logic is critical and complex enough to warrant
  its own unit tests that can be performed in isolation.
\end{itemize}

A lambda is often appropriate for self-contained, cohesive blocks of logic, even if they are several
lines long. It is when the logic loses this cohesion or makes the surrounding code difficult to
follow that a refactor is warranted.

For example, a stream mapping operation that involves several steps of transformation and validation
is a prime candidate for extraction.

\begin{snippet}
// AVOID!
List<Integer> results = strings.stream().map(s -> {
    String trimmed = s.trim();
    if (trimmed.length() < 2) {
        return 0;
    }
    try {
        int value = Integer.parseInt(trimmed.substring(0, 2));
        // Further complex business logic...
        return value * 2;
    } catch (NumberFormatException _) {
        LOGGER.warning(() -> "Invalid number format: " + s);
        return -1; // return a default value on parsing failure
    }
}).toList();


// RIGHT: LOGIC EXTRACTED INTO A PRIVATE METHOD.
List<Integer> results = strings.stream().map(this::parseAndTransformString).toList();

...

private Integer parseAndTransformString(String s) {
    String trimmed = s.trim();
    if (trimmed.length() < 2) {
        return 0;
    }
    try {
        int value = Integer.parseInt(trimmed.substring(0, 2));
        // Further complex business logic...
        return value * 2;
    } catch (NumberFormatException _) {
        LOGGER.warning(() -> "Invalid number format: " + s);
        return -1; // Return a default value on parsing failure
    }
}
\end{snippet}

\subsubsection{Method References}

A lambda expression that does nothing more than call an existing method by name must be replaced
with a method reference (\texttt{::}).

\begin{snippet}
// AVOID!
stream.forEach(s -> System.out.println(s));
stream.map(s -> s.trim());
stream.map(s -> new BigDecimal(s));

// RIGHT
stream.forEach(System.out::println);
stream.map(String::trim);
stream.map(BigDecimal::new);
\end{snippet}

\subsubsection{Unused Lambda Parameters}

When a lambda parameter is not used in the body of the lambda, its name should be an unnamed
variable (\texttt{\_}).

\begin{snippet}
map.forEach((key, value) -> System.out.println(value));  // AVOID!

map.forEach((_, value) -> System.out.println(value));    // OK
\end{snippet}

\subsection{Streams API}

The Streams API provides a powerful mechanism for processing sequences of elements. When used
appropriately, streams can make code more declarative and readable.

Streams should be used for data processing operations where the sequence of operations is clear and
concise. Avoid creating long or overly complex stream pipelines, as they can become difficult to
debug and understand.

A good rule of thumb is to refactor any stream pipeline that consists of more than three or four
intermediate operations. This can be done by extracting the logic into a helper method or by
breaking the pipeline into multiple, distinct statements.

\begin{snippet}
// AVOID!
List<String> results = names.stream()
        .filter(s -> s.length() > 3)
        .map(s -> s.substring(0, 2))
        .sorted()
        .distinct()
        .filter(s -> !s.equals("NO"))
        .map(s -> s + "!")
        .collect(Collectors.toList());

// RIGHT
List<String> results = names.stream()
        .filter(s -> s.startsWith("A"))
        .map(String::toUpperCase)
        .toList();
\end{snippet}

\subsection{Text Blocks}

Text blocks should be used for multi-line string literals. They improve readability for strings
spanning multiple lines, such as embedded snippets of SQL, JSON, or HTML. Text blocks eliminate the
visual clutter of most escape sequences and string concatenations.

A text block is opened with three double-quote characters followed by a line terminator. The content
of the string begins on the next line.

\begin{snippet}
String json = """
        {
          "user": "test",
          "active": true
        }
        """;
\end{snippet}

This form is preferred over the equivalent string literal concatenation:

\begin{snippet}
String json = "{\n"
        + "  \"user\": \"test\",\n"
        + "  \"active\": true\n"
        + "}\n";
\end{snippet}

\subsection{Pattern Matching}

Pattern matching streamlines common programming tasks by combining type testing, casting, and
variable declaration into a single, concise expression. Its use is recommended to improve code
readability and reduce boilerplate.

\subsubsection{Pattern Matching for \texttt{instanceof}}

A pattern matching \texttt{instanceof} check should be used to test, cast, and bind a variable in a
single expression. This form is preferred over a separate \texttt{instanceof} check followed by an
explicit cast.

\begin{snippet}
// AVOID!
if (obj instanceof String) {
    String s = (String) obj;
    System.out.println("String length: " + s.length());
}

// RIGHT
if (obj instanceof String s) {
    System.out.println("String length: " + s.length());
}
\end{snippet}

\subsubsection{Pattern Matching for \texttt{switch}}

A pattern matching \texttt{switch} should have the following form to test against multiple types
and patterns. This is the preferred modern alternative to long \texttt{if-else-if} chains.

When switching on a variable that can be \texttt{null}, prefer a \texttt{case null} label over an
explicit \texttt{null} check guarding the switch statement. Conditional logic can be applied to a
pattern using a \texttt{when} clause. An unnamed pattern (\texttt{\_}) should be used for components
of a record or pattern that are intentionally ignored.

\begin{snippet}
// AVOID!
if (obj == null) {
    System.out.println("Object is null");
} else if (obj instanceof String) {
    String s = (String) obj;
    System.out.println("String: " + s.toUpperCase());
} else if (obj instanceof Integer) {
    Integer i = (Integer) obj;
    if (i > 100) {
        System.out.println("Large integer");
    } else {
        /* Handle other integers. */
    }
} else if (obj instanceof Point) {
    Point p = (Point) obj;
    System.out.println("A point with x-coordinate: " + p.x());
} else {
    /* Handle all other types. */
}

// RIGHT
switch (obj) {
    case null -> System.out.println("Object is null");
    case String s -> System.out.println("String: " + s.toUpperCase());
    case Integer i when i > 100 -> System.out.println("Large integer");
    case Integer i -> { /* Handle other integers */ }
    case Point(int x, _) -> System.out.println("A point with x-coordinate: " + x);
    default -> { /* Handle all other types */ }
}
\end{snippet}

\begin{note}
When used with \texttt{sealed} classes and interfaces, a pattern matching \texttt{switch} can be
checked for exhaustiveness by the compiler, eliminating the need for a default case and preventing
a common source of bugs.
\end{note}

\section{Code Example}

\subsection{Java Source File Example}
\label{codeexample}

The following example shows how to format a Java source file containing a single public class.
Interfaces are formatted similarly. For more information, see
\hyperref[classesinterfaces]{"\nameref*{classesinterfaces}" on page~\pageref*{classesinterfaces}}
and \hyperref[doccomments]{"\nameref*{doccomments}" on page~\pageref*{doccomments}}.

\begin{lstlisting}[
  language=Java,
  morekeywords={enum, record},
  commentstyle=\color{green!60!black}\itshape,
  stringstyle=\color{red!40!black},
  keywordstyle=\color{blue}\bfseries,
  numbers=left,
  numberstyle=\scriptsize\color{gray},
  numbersep=5pt,
  linewidth=\dimexpr\textwidth\relax]
/*
 * Copyright (c) 2025 Larry Kramer
 *
 * Permission is hereby granted, free of charge, to any person obtaining a copy
 * of this software and associated documentation files (the "Software"), to deal
 * in the Software without restriction, including without limitation the rights
 * to use, copy, modify, merge, publish, distribute, sublicense, and/or sell
 * copies of the Software, and to permit persons to whom the Software is
 * furnished to do so, subject to the following conditions:
 *
 * The above copyright notice and this permission notice shall be included in
 * all copies or substantial portions of the Software.
 *
 * THE SOFTWARE IS PROVIDED "AS IS", WITHOUT WARRANTY OF ANY KIND, EXPRESS OR
 * IMPLIED, INCLUDING BUT NOT LIMITED TO THE WARRANTIES OF MERCHANTABILITY,
 * FITNESS FOR A PARTICULAR PURPOSE AND NONINFRINGEMENT. IN NO EVENT SHALL THE
 * AUTHORS OR COPYRIGHT HOLDERS BE LIABLE FOR ANY CLAIM, DAMAGES OR OTHER
 * LIABILITY, WHETHER IN AN ACTION OF CONTRACT, TORT OR OTHERWISE, ARISING FROM,
 * OUT OF OR IN CONNECTION WITH THE SOFTWARE OR THE USE OR OTHER DEALINGS IN THE
 * SOFTWARE.
 */

package net.larrykramer.somepackage;

import module java.logging;

import java.io.IOException;
import java.util.ArrayList;
import java.util.List;

import static java.lang.Math.max;

/**
 * Manages HTTP network connections for operations.
 * <p>
 * This class demonstrates all coding standards, including formatting, naming, comments, imports,
 * control structures, and more.
 */
public class NetworkManager {
    /** The maximum number of concurrent connections allowed. */
    public static final int MAX_CONNECTIONS = 10;

    private static final Logger LOGGER = Logger.getLogger(NetworkManager.class.getName());

    private static final int DEFAULT_TIMEOUT_SECONDS = 30;

    private List<String> activeConnections;
    private int maxConnections;

    /**
     * Constructs a new NetworkManager.
     */
    public NetworkManager() {
        this(MAX_CONNECTIONS);
    }

    /**
     * Constructs a new NetworkManager with a given number of connections.
     *
     * @param count the maximum number of connections this manager can handle
     * @throws IllegalArgumentException if {@code count < 0} or {@code count > MAX_CONNECTIONS}
     */
    public NetworkManager(int count) {
        if (count < 0) {
            throw new IllegalArgumentException("count cannot be negative");
        }
        if (count > MAX_CONNECTIONS) {
            throw new IllegalArgumentException("count cannot be greater than " + MAX_CONNECTIONS);
        }
        this.maxConnections = count;
        this.activeConnections = new ArrayList<>();
    }

    /**
     * Opens a new network connection.
     *
     * @param url the URL to connect to
     * @throws IOException if the URL is invalid or a connection cannot be established
     */
    public void openConnection(String url) throws IOException {
        if (activeConnections.size() >= maxConnections) {
            throw new IOException("Maximum connections reached.");
        }

        // Simulate opening a network connection
        activeConnections.add(url);
    }

    /**
     * Closes an existing network connection.
     *
     * @param url the URL of the connection to close
     */
    public void closeConnection(String url) {
        if (activeConnections.contains(url)) {
            activeConnections.remove(url);
        }
    }

    /**
     * Returns the number of active connections.
     *
     * @return the current number of active connections
     */
    public int getActiveConnectionCount() {
        return activeConnections.size();
    }

    /**
     * Performs a network reset based on status.
     *
     * @param status the current network status
     */
    public void handleNetworkStatus(NetworkStatus status) {
        // The example below uses a traditional switch to demonstrate /* falls through */
        // comments, but new code should prefer switch expressions where possible.
        switch (status) {
            case CONNECTED:
                // No action needed
                break;
            case DISCONNECTED:
                resetAllConnections();
                break;
            case ERROR:
                resetAllConnections();
                /* falls through */
            default:
                logStatus("Unknown status, performing reset.");
                resetAllConnections();
                break;
        }
    }

    private void resetAllConnections() {
        // Clear all active connections safely
        activeConnections.clear();
    }

    private void logStatus(String message) {
        LOGGER.info(message);
    }

    /**
     * Calculates the optimal timeout based on load.
     *
     * @param connectionCount the number of active connections
     * @return the calculated timeout in seconds
     */
    public static int calculateTimeout(int connectionCount) {
        // Use a named constant instead of a magic number
        final int timeoutFactor = 2;
        return max(DEFAULT_TIMEOUT_SECONDS, connectionCount * timeoutFactor);
    }

    /** Represents possible network statuses. */
    public enum NetworkStatus {
        CONNECTED,
        DISCONNECTED,
        ERROR
    }

    /**
     * This record encapsulates the connection details.
     */
    record ConnectionDetail(String url, long timestamp) {
    }
}
\end{lstlisting}

\appendix

\section{Intentional Divergences and References}

\subsection{Divergences from Sun/Oracle \emph{Java Code Conventions}}

\begin{itemize}
\item Line length is 100 characters (Sun/Oracle recommended 80).
\item Indentation is spaces only; tabs are disallowed except in very specific circumstances
  (Sun/Oracle allowed tabs, recommending they be set to 8 spaces).
\item Import ordering and wildcard thresholds are explicitly defined (Sun/Oracle did not prescribe
  these).
\item One blank line is used between sections of a source file (Sun/Oracle recommended two).
\item Variables are declared close to their first use (Sun/Oracle required declarations at the
  beginning of a block).
\item Avoid redundant parentheses when precedence is clear (Sun/Oracle encouraged their liberal
  use).
\end{itemize}

\subsection{References}

\begin{itemize}
\item \href{https://www.oracle.com/docs/tech/java/codeconventions.pdf}{Sun/Oracle \emph{Java Code Conventions} (1997 Edition)}
\item \href{https://www.oracle.com/technical-resources/articles/java/javadoc-tool.html}{Oracle: How to Write Doc Comments for the Javadoc Tool}
\item \href{https://openjdk.org/jeps/126}{JEP 126: Lambda Expressions \& Virtual Extension Methods}
\item \href{https://openjdk.org/jeps/286}{JEP 286: Local-Variable Type Inference}
\item \href{https://openjdk.org/jeps/361}{JEP 361: Switch Expressions}
\item \href{https://openjdk.org/jeps/378}{JEP 378: Text Blocks}
\item \href{https://openjdk.org/jeps/394}{JEP 394: Pattern Matching for instanceof}
\item \href{https://openjdk.org/jeps/395}{JEP 395: Records}
\item \href{https://openjdk.org/jeps/441}{JEP 441: Pattern Matching for switch}
\item \href{https://openjdk.org/jeps/456}{JEP 456: Unnamed Variables \& Patterns}
\item \href{https://openjdk.org/jeps/511}{JEP 511: Module Import Declarations}
\item \href{https://docs.oracle.com/en/java/javase/21/docs/api/java.base/java/util/stream/package-summary.html}{Oracle Java 21 API: Stream API}
\end{itemize}

\end{document}
